\chapter{相关技术概述}\label{chap:related-work}

\section{系统架构与服务组织技术}

本系统采用前后端分离的总体实现思路。前端以微信小程序作为统一交互入口，负责界面展示、状态维护和用户操作响应；后端则通过独立服务承担用户管理、拼车处理、论坛交流、聊天通信和交通信息查询等业务功能。该架构将界面逻辑与业务逻辑分离，有利于降低模块之间的耦合程度，并提高后续维护和扩展的便利性。

在后端组织方式上，系统采用按业务能力拆分的微服务结构。不同服务围绕相对独立的业务边界进行部署和演化，并通过 REST 接口完成数据交换与流程协同。对于本系统而言，这种方式能够较好地适应校园综合交通场景中业务种类多、数据类型差异大以及后续功能持续扩展的特点。

\section{微信小程序前端技术}

微信小程序是本系统面向用户的主要访问载体。其运行环境具备轻量化、免安装和与校园用户日常使用习惯高度一致等特点，适合作为校内交通服务的统一入口。系统前端基于小程序原生框架构建，通过页面配置、组件封装与事件绑定实现论坛、拼车、交通、消息和个人中心等主要功能页面。

在工程实现中，小程序前端主要依赖 WXML、WXSS 和 JavaScript 完成界面组织与交互逻辑处理。系统通过页面生命周期函数进行数据初始化，通过 \texttt{setData} 驱动视图更新，并通过自定义组件复用高频界面单元。与此同时，前端还结合微信提供的授权登录、图片选择、订阅消息和本地缓存等能力，完成用户登录、媒体上传和业务提醒等功能支持。

\section{Spring Boot 后端开发技术}

后端服务基于 Spring Boot 框架实现。Spring Boot 具备自动配置、依赖整合和快速构建 Web 服务等特点，能够显著降低基础工程搭建成本，使开发重点集中在业务逻辑实现上。对于本系统而言，Spring Boot 为控制层、业务层和数据访问层的清晰组织提供了良好支撑，也便于不同微服务以统一风格进行开发。

在服务实现中，系统使用注解方式组织控制器、业务服务和数据访问对象，并通过统一响应结构和集中异常处理提升接口的一致性。对于跨服务调用，系统通过封装客户端组件完成 HTTP 交互，使拼车服务能够访问用户服务和模型服务，论坛服务能够获取用户资料，从而实现多模块协同工作。该技术路线兼顾了工程实现效率与系统结构清晰度。

\section{数据库与文档存储技术}

为适应不同业务数据的特征，本系统采用 MySQL 与 MongoDB 组合的混合存储方案。MySQL 主要用于存储用户、拼车订单、拼车申请和拼车反馈等结构稳定且具有较强事务一致性要求的数据。这类数据通常伴随明确的主键关系、状态流转和数量约束，适合通过关系型数据库进行管理。

MongoDB 主要用于存储帖子、评论、聊天会话、聊天消息、文件记录及交通时刻表等半结构化或嵌套型数据。这类数据在字段组织上更灵活，往往包含数组、嵌套对象和可变结构，使用文档数据库能够减少复杂表连接，提高内容型数据和消息型数据的读写适应性。通过两类数据库的协同使用，系统在一致性、灵活性和扩展性之间取得了较好的平衡。

\section{本章小结}

本章概述了系统实现所依赖的关键技术，包括前后端分离架构、微信小程序开发、Spring Boot 微服务以及 MySQL 与 MongoDB 的混合存储方案。这些技术共同构成了平台实现的基础，也为后续需求分析、系统设计和功能实现提供了技术支撑。
